% HD-S32. GPT-6 Astra Pro; inherited effort not exposed; /root/sixfolds.
% 21 November 2026. Exact root matrices, integral pairing and actual factor.
\subsection{HD-S32: all natural \(G_{29}\) and \(G_{31}\) Weil sixfolds have determinant \([-1]\) or \([-3]\)}
\label{sec:gaussian-rank-four-scope}

Let \(G\) be either of the explicit rank-four Gaussian reflection
groups below. For every connected sphere cover with full monodromy
\(G\), its natural primitive Jacobian factor of signature \((3,3)\),
equipped with the restricted theta polarization, has ordinary
Hermitian determinant class \([-1]\) or \([-3]\). This conclusion
holds for every branch profile. The proof retains the coefficient
lattice, local boundary maps, and the map to the actual polarized
factor. It also describes precisely the allowed scalar extensions
and changes of polarization in a rank-nine family.

\paragraph{1. Published roots and a self-dual coefficient lattice.}
Write \(u_1,\ldots,u_4\) for the standard orthonormal basis of
\(K^4\), where \(K=\mathbb Q(i)\). Put
\[
\begin{split}
 v_1&=\tfrac{1+i}{2}(u_2+u_4),\qquad
 v_2=\tfrac{1+i}{2}(u_3-u_2),\\
 v_3&=\tfrac12(-u_1+iu_2-u_3+iu_4),\qquad v_4=u_1,\\
 v_5&=\tfrac12(u_1+u_2+u_3+u_4),\qquad
 r_j=I_4-2v_jv_j^*.
\end{split}
\]
All five vectors have norm one. Cohen, \emph{Finite complex
reflection groups}, Ann. Sci. ENS 9 (1976), (4.15), printed p.408,
gives the first four as the root graph \(N_4\). The theorem
(4.15)(iii), printed p.409, adjoins the direction
\(u_1+u_2+u_3+u_4\) to obtain \(EW(N_4)\). Its table (4.16),
printed p.412, identifies these groups with \(G_{29}\) and
\(G_{31}\), respectively. Thus
\[
 G_{29}=\langle r_1,r_2,r_3,r_4\rangle,
 \qquad G_{31}=\langle r_1,r_2,r_3,r_4,r_5\rangle.
\]
The primary source is
\url{https://www.numdam.org/item/ASENS_1976_4_9_3_379_0.pdf}.
Both root pages and the table have been inspected in the original
scan, including the standard-basis notation in the extension.

Directly from the displayed vectors, every \(r_j\) has its entries
in \(\tfrac12\mathbb Z[i]\), satisfies \(r_j^2=I_4\), and
preserves the Hermitian Gram matrix \(I_4\). Consequently
\[
 R=\mathbb Z[i,1/30],\qquad L=R^4,\qquad h_0(x,y)=\bar x^t y
 \tag{S32.1}
\]
is an explicit invariant self-dual coefficient lattice. In particular,
the support of its discriminant has been determined from the actual
Gram matrix, rather than inferred from the group order.

The exact computation \texttt{gaussian\_reflection\_scope\_certificate.g}
starts from these five matrices. It closes the four standard row
vectors under their right action. The resulting row orbits have
sizes 160 and 240. Their permutation actions are faithful because
they contain a basis, and the matrix attached to a permutation is
recovered from the four basis-row images. With GAP's right-action
convention this is the genuine representation
\(p(MN)=p(M)p(N)\); the script checks the generator products.
Fresh group closure and conjugacy-class computation give
\[
\begin{array}{c|r|r|r|r}
G&|G|&\#\text{classes}&\#\text{reflections}&|G|/4\\\hline
G_{29}&7680=2^9\cdot3\cdot5&37&40&1920\\
G_{31}&46080=2^{10}\cdot3^2\cdot5&59&60&11520
\end{array}
\]
Every reflection has order two. Both groups contain the actual scalar
matrix \(z=iI_4\). The class sum of \(|\operatorname{tr}\rho(g)|^2\)
is \(|G|\), proving absolute irreducibility. The matrices are defined
over \(K\), and the value \(\operatorname{tr}(z)=4i\) shows that
their character field is exactly \(K\). The complete class data are
saved in \texttt{g29\_scope\_certificate.json} and
\texttt{g31\_scope\_certificate.json}; no group, class or character
table is imported by the computation.

\paragraph{2. The exact local splittings.}
For any local monodromy \(g\), let \(o=\operatorname{ord}(g)\).
Its order divides the displayed group order and is a unit in \(R\).
Define the integral \(R\)-endomorphisms
\[
 e_g=\frac1o\sum_{j=0}^{o-1}g^j,\qquad
 B_g=\frac1o\sum_{j=0}^{o-1}(o-1-j)g^j.
 \tag{S32.2}
\]
A direct telescoping computation gives
\[
 e_g^2=e_g,\quad ge_g=e_g,\quad
 (1-g)B_g=1-e_g.
\]
It follows that \(L^g=e_gL\), and the map induced by \(e_g\)
is the isomorphism \(L/(g-1)L\simeq e_gL\). Both boundary
modules are free direct summands of the coefficient lattice.

The exact class computation verifies these identities for every
representative. It also computes the moving factor
\(f_g(T)=\det(T-g)/(T-1)^{\dim L^g}\) and finds
\[
\begin{array}{c|l}
G_{29}&N(f_g(1))\in\{1,4,9,16,25,36,64,256\}\\
G_{31}&N(f_g(1))\in\{1,4,9,16,25,36,64,81,256\}.
\end{array}
\]
Thus one can alternatively use the actual resultant projector
\(f_g(g)/f_g(1)\) and the Bezout map of HD-S28. The elementary
averages (S32.2) already give all required local splittings for
every possible branch profile.

\paragraph{3. The integral parabolic pairing.}
Let \(\Sigma\) be the sphere with the branch discs removed, with
coefficient system \(L\). Full monodromy contains \(z=iI_4\),
and \(i-1\) is a unit of \(R\). For any monodromy word \(g\),
\((g-1)L\) is contained in the sum of the generator images
\((g_j-1)L\). Therefore the absolute zeroth homology of both
\(L\) and its dual is zero. Apply the bordered-surface argument
proved in HD-S27, paragraph 2, over the PID \(R\), with the local
splittings (S32.2).

To specify the resulting integral map, absolute first homology is
the kernel of the cellular map for a graph retract of \(\Sigma\).
Poincar\'e--Lefschetz cap product and the universal coefficient
sequence identify it perfectly with the dual of relative first
homology for the dual coefficient system. Both modules are free.
The image
\[
 P(L)=\operatorname{im}\bigl(H_1(\Sigma,L)
              \longrightarrow H_1(\Sigma,\partial\Sigma;L)\bigr)
\]
has free cokernel \(H_0(\partial\Sigma,L)\), by the pair sequence
and the vanishing of absolute \(H_0\). Hence it is primitive.
The absolute-to-relative maps are adjoint under the two perfect
cross-pairings. Extending any functional from the primitive image
to the relative module gives the actual isomorphism from the dual
parabolic image to \(P(L)^\vee\); adjointness identifies its
kernel before passage to the image. This proves a perfect integral
pairing, including surjectivity. The Mayer--Vietoris maps for the
branch discs identify this same image with
\(H^1(\mathbb P^1,j_*L)\), because the local averages make the
disc extensions invariant summands with no higher cohomology.

Using the self-dual form (S32.1), this is a perfect skew-Hermitian
parabolic cup form \(S\). In primitive dimension six, its
\(R\)-rank is six. The Hermitian form \(H=iS\), with the
Gaussian action convention of HD-S23, is again perfect. A primary
cap-product source for these maps is Dettweiler--Wewers,
Proposition 2.1, printed pp.153--154; the extension from
nondegeneracy to integral perfection here is the explicit primitive
image argument just given.

Since \(H\) is Hermitian, its determinant is rational. Perfection
puts it in \(R^\times\); signature \((3,3)\) makes it negative.
Thus
\[
 \det H=-2^a3^b5^c\quad(a,b,c\in\mathbb Z).
\]
The identities \(2=N(1+i)\), \(5=N(2+i)\), and
\(3^2=N(3)\) now give exactly the two possible ordinary classes
\([-1]\) and \([-3]\).

\paragraph{4. The literal primitive factor and its polarization.}
Take \(v=(1+i)u_1\), so \(h_0(v,v)=2\), and set
\[
 e=\frac4{|G|}\sum_{g\in G}
       \operatorname{Re}h_0(v,\rho(g^{-1})v)\,g,
 \qquad D=|G|/4.
 \tag{S32.3}
\]
The matrix-coefficient orthogonality calculation of HD-S23 proves
\(e^2=e=e^*\) and \(e\mathbb Q[G]e=K\): on the natural
complex block it is the orthogonal projector onto the line spanned
by \(v\), and on the conjugate block it is its conjugate; all
other irreducible blocks vanish. The row-orbit computation also
checks that every group matrix has entries in
\(\tfrac12\mathbb Z[i]\). Thus the coefficients
\(\operatorname{Re}h_0(v,\rho(g^{-1})v)
 =2\operatorname{Re}\rho(g^{-1})_{11}\) are integers, and
\(De\in\mathbb Z[G]\).

For the left deck action \(L_g\), use the actual covariant
Jacobian endomorphism \(a_g=L_{g,*}\). Define
\(B=\operatorname{im}(De:J(C)\to J(C))\) and
\(i_B=a_z^{-1}|_B\). On curve cohomology put
\(T_g=(L_g^{-1})^*\). The Abel map identifies \(a_g^*\) with
\(T_{g^{-1}}\), so \(i_B^*=T_z\). These are the same
left-deck and duality conventions used in HD-S23.

The coefficient map is the function
\(j_x(g)=h_0(v,\rho(g)x)/2\). Its complex Schur identity and
the real map \(\Phi=\operatorname{PD}(j+\bar j)\) give
\[
 \sum_g\bar j_x(g)j_y(g)=\frac{|G|}{8}h_0(x,y),
 \qquad h_B(\Phi x,\Phi y)=D H(x,y).
 \tag{S32.4}
\]
Indeed the real alternating cup form receives the two conjugate
cross terms, so the complex scalar \(|G|/8\) becomes \(|G|/4\).
The same-block products vanish: the central scalar acts on them by
\(-1\), whereas it acts trivially on their invariant cup value.
The Poincar\'e-duality convention then gives precisely the
Hermitian identity displayed in (S32.4), and
\(\Phi(ix)=a_z^{-1}{}_\ast\Phi(x)\). This is an actual
\(K\)-linear rational polarized isomorphism to the primitive
factor, not an identification by dimensions. Its determinant
multiplier \(D^6=N(D^3)\) is a norm. Consequently the two-class
conclusion holds for the actual restricted-theta factor.

Changing the rank-one coefficient line uses the explicit
matrix-coefficient intertwiner; its Hermitian scalar is rational
and positive, hence again has a sixth-power norm determinant.
Any specified \(K\)-linear isometry or rational similitude to a
different lattice has the same conclusion. For a rank-nine family,
a fixed isogeny and another compatible flat polarization also
preserve it: pull the second form back along the isogeny, and apply
the explicit common-centralizer calculation of HD-S17 on the open
unitary period domain. Its self-adjoint scalar is in
\(K\cap\mathbb R=\mathbb Q\), so its rank-six determinant
multiplier is a norm. This last assertion uses the actual open
period family; it makes no assertion about an arbitrary additional
polarization at an isolated special point.

\paragraph{5. Source scope and checks.}
A finite scalar extension inside \(GL_4(K)\) can only adjoin
roots of unity of \(K\), namely \(\mu_4\). Both groups already
contain \(iI_4\), so such an extension changes neither group.
The inclusion \(G_{29}\subset G_{31}\) is the literal inclusion
of their displayed matrix generators. No assertion about other
nonscalar overgroups is used.

For orientation, the parabolic dimension formula for a sixfold is
\(\sum_j\operatorname{rank}(1-g_j)=14\). A rank-nine period
map requires at least twelve branch values, so
\(\sum_j(\operatorname{rank}(1-g_j)-1)\le2\). The determinant
proof above does not depend on this restriction and covers all
natural sixfold profiles of the two groups. It supplies no new
dominance or cycle assertion.

The whole argument was read by its author, and the exact matrix
certificate was executed from the published roots. Group order,
all conjugacy classes, every local projector, the coefficient Gram
matrix, and the integral coefficients of (S32.3) are separate
checked inputs. No universal determinant conclusion is inferred
from a sample of tuples. The event log records these checks and
any later independent review distinctly.
